\documentclass[12pt,a4paper]{article}
\usepackage{amsmath,amssymb,amsthm}
\usepackage{geometry}
\usepackage{hyperref}
\usepackage{url}
\usepackage{booktabs}
\usepackage{xcolor}
\usepackage{enumitem}
\usepackage{microtype}
\geometry{margin=2.5cm}
\hypersetup{colorlinks=true,linkcolor=blue,citecolor=blue,urlcolor=blue}
\newtheorem{theorem}{Theorem}
\newtheorem{proposition}[theorem]{Proposition}
\newtheorem{lemma}[theorem]{Lemma}
\newtheorem{conjecture}[theorem]{Conjecture}
\theoremstyle{definition}
\newtheorem{definition}[theorem]{Definition}
\newtheorem{remark}[theorem]{Remark}
\newcommand{\ex}{\mathrm{ex}}

\title{\textbf{New Lower Bounds for $C_4$-Free Subgraphs\\
of the Hypercubes $Q_6$, $Q_7$, and $Q_8$:\\
Constructions, Structure, and Computational Method}}
\author{Minamo Minamoto\\[4pt]
\small Independent\\
\small \texttt{minamominamoto4f5683f6@gmail.com}}
\date{March 2026}

\begin{document}
\maketitle

\begin{abstract}
We establish new lower bounds $\ex(Q_7,C_4)\ge304$ and $\ex(Q_8,C_4)\ge680$
for the maximum number of edges in a $C_4$-free subgraph of the 7- and
8-dimensional hypercubes, and give a modern computational reproduction of
$\ex(Q_6,C_4)=132$.
All bounds are witnessed by explicit constructions supplied as edge lists with
machine-verifiable certificates: a complete, deterministic enumeration of all
four-cycles (240 for $Q_6$, 672 for $Q_7$, 1\,792 for $Q_8$) establishes
$C_4$-freeness, and each certificate is reproducible by an independent,
dependency-free verifier.

For $Q_7$ we identify $19\,866$ distinct $C_4$-free subgraphs on 304 edges;
their \emph{dimension profiles} fall into exactly 20 types.  All $19\,866$
solutions share a rigid structural core: degree sequence
$\{4^{32},5^{96}\}$, spectral radius $\lambda_1\approx4.787$, and local
maximality.  Pairwise Hamming distances range from 36 to 260.
Whether these solutions exhaust all 304-edge $C_4$-free subgraphs of
$Q_7$ remains open.

For $Q_8$ we analyse the local structure of the 680-edge construction:
every non-edge of the construction creates at least one $C_4$, and
1\,076 independent searches at 681 edges did not achieve zero violations.
These observations constitute computational evidence, not a proof,
of the conjectured equality $\ex(Q_8,C_4)=680$.

The constructions are found by a two-phase simulated annealing algorithm
with Aut$(Q_n)$-based diversification, described in full detail.
For $Q_6$ we additionally provide an ILP-based proof that
$\ex(Q_6,C_4)\le132$, reproducing the exact value computationally.

Edge lists, ILP files, and source code are publicly available at\\
\url{https://github.com/minamominamoto/c4free-hypercube}.
\end{abstract}

\medskip
\noindent\textbf{MSC 2020:} 05C35, 05C62, 05C50.\\
\textbf{Keywords:} hypercube, $C_4$-free, extremal graph theory,
simulated annealing, integer linear programming, dimension profile.

\tableofcontents

%% ============================================================
\section{Introduction}
%% ============================================================

The $n$-dimensional hypercube $Q_n$ has vertex set $\{0,1\}^n$, two vertices
adjacent iff they differ in exactly one coordinate; it has $2^n$ vertices and
$n\cdot2^{n-1}$ edges.  The problem of determining
$\ex(Q_n,C_4)=\max\{|E(G)|:G\subseteq Q_n,\,G\text{ is }C_4\text{-free}\}$
was raised by Erd\H{o}s~\cite{Erd84} and remains open in general.

The best known asymptotic bounds are
\begin{equation}\label{eq:bounds}
\tfrac{1}{2}(n+\sqrt{n})\cdot2^{n-1}\;\le\;\ex(Q_n,C_4)\;\le\;0.6032\cdot n\cdot2^{n-1},
\end{equation}
where the lower bound (valid for $n=4^r$) is due to Brass, Harborth, and
Nienborg~\cite{BHN95}, and the upper bound to Baber~\cite{Bab12} via the flag
algebra method, improving the earlier bound of Thomason and Wagner~\cite{TW09}.
Balogh et al.~\cite{BHLL12} give 0.6068 via a related method.
For general $n\ge9$, the weaker bound
$\tfrac{1}{2}(n+0.9\sqrt{n})\cdot2^{n-1}$ applies~\cite{BHN95}.
Exact values have been determined for small $n$; Harborth and Nienborg~\cite{HN94}
proved $\ex(Q_6,C_4)=132$.

Our main results are:

\begin{theorem}\label{thm:main}
\begin{enumerate}[label=(\roman*)]
  \item $\ex(Q_7,C_4)\ge304$.
  \item $\ex(Q_8,C_4)\ge680$.
\end{enumerate}
\end{theorem}

Both bounds are realised by explicit constructions (supplementary files
\texttt{q7\_edges\_304.jsonl} and \texttt{q8\_edges\_680.jsonl}).

\begin{table}[ht]
\centering
\caption{Known values and new lower bounds of $\ex(Q_n,C_4)$.
  The BHN column gives the Brass--Harborth--Nienborg general estimate.}
\label{tab:values}\medskip
\begin{tabular}{ccccl}
\toprule
$n$ & $|E(Q_n)|$ & $\ex(Q_n,C_4)$ & BHN est. & Source\\
\midrule
4 &   32 & 24  & $\approx25.6$ & Chung~\cite{Chu92}\\
5 &   80 & 56  & $\approx51.2$ & Emamy-K et al.~\cite{EMR92}\\
6 &  192 & 132 & $\approx96.0$ & Harborth--Nienborg~\cite{HN94} (reproduced)\\
7 &  448 & $\ge\mathbf{304}$ & $\approx300.2$ & This work\\
8 & 1024 & $\ge\mathbf{680}$ & $\approx674.9$ & This work\\
\bottomrule
\end{tabular}
\end{table}

The paper is organised as follows.
Section~\ref{sec:verify} describes the $C_4$ enumeration and verification
procedure used throughout.
Sections~\ref{sec:q7}--\ref{sec:q8} present the $Q_7$ and $Q_8$ constructions.
Section~\ref{sec:structure} classifies the $19\,866$ $Q_7$ solutions found
by their dimension profiles, yielding exactly 20 types.
Section~\ref{sec:method} describes the two-phase simulated annealing algorithm.
Section~\ref{sec:q6} reproduces the exact value $\ex(Q_6,C_4)=132$ via ILP.
Section~\ref{sec:open} lists open problems.

%% ============================================================
\section{$C_4$ Enumeration and Verification}\label{sec:verify}
%% ============================================================

\begin{lemma}[$C_4$ enumeration in $Q_n$]\label{lem:c4}
Every four-cycle of $Q_n$ is uniquely determined by a pair of coordinate
directions $0\le i<j\le n-1$ and a base vertex $b\in\{0,1\}^n$ with
$b_i=b_j=0$.  Its four vertices are
\[
  b,\quad b\oplus e_i,\quad b\oplus e_j,\quad b\oplus e_i\oplus e_j,
\]
and the total number of four-cycles in $Q_n$ is $\binom{n}{2}\cdot2^{n-2}$.
For $n=6$: $240$; for $n=7$: $672$; for $n=8$: $1\,792$.
\end{lemma}

\noindent\textbf{Verification algorithm.}\quad
Given a candidate edge set $E\subseteq E(Q_n)$ stored as a hash set,
the following procedure performs a complete, deterministic $C_4$-freeness check:
\begin{quote}\small
\textbf{for each} pair $(i,j)$ with $0\le i<j\le n-1$:\\
\quad\textbf{for each} $b\in\{0,1\}^n$ with $b_i=b_j=0$:\\
\quad\quad let $v_1=b,\;v_2=b\oplus e_i,\;v_3=b\oplus e_j,\;v_4=b\oplus e_i\oplus e_j$\\
\quad\quad \textbf{if} $\{v_1v_2,\,v_1v_3,\,v_2v_4,\,v_3v_4\}\subseteq E$:
  \textbf{return} \textsc{violation}\\
\textbf{return} \textsc{c4-free}
\end{quote}
We applied this algorithm to all constructions and all $19\,866$ solutions
of $Q_7$; every call returned \textsc{c4-free}.

%% ============================================================
\section{The $Q_7$ Construction}\label{sec:q7}
%% ============================================================

We constructed a $C_4$-free subgraph of $Q_7$ on 304 edges by
penalty-based simulated annealing (Section~\ref{sec:method}),
then certified it by Lemma~\ref{lem:c4}'s algorithm.

\begin{proposition}\label{prop:q7}
The 304-edge construction has the following properties:
\begin{enumerate}[label=(\roman*)]
  \item Degree sequence: 32 vertices of degree 4 and 96 of degree 5
    (degree sum $608=2\cdot304$).
  \item Bipartite: $\lambda_1\approx4.787$, $\lambda_n\approx-4.787$,
    $\lambda_1+\lambda_n=0$ to machine precision.
  \item $\mathrm{Tr}(A^4)=5216$, attaining the $C_4$-free trace equality
    $\sum_i\deg(i)^2+\sum_{uv\in E}(\deg u+\deg v)-2|E|$.
  \item Local maximality: every non-edge of $Q_7$ creates a $C_4$ when added.
\end{enumerate}
\end{proposition}

Through $19\,866$ independent runs (Section~\ref{sec:landscape})
we found $19\,866$ $C_4$-free subgraphs of $Q_7$ on 304 edges with
pairwise distinct edge sets (not quotiented by $\mathrm{Aut}(Q_7)$),
all sharing the structural properties stated in Proposition~\ref{prop:rigid}.

%% ============================================================
\section{The $Q_8$ Construction}\label{sec:q8}
%% ============================================================

Simulated annealing produced a $C_4$-free subgraph of $Q_8$ on 680 edges,
certified by exhaustive enumeration of all $1\,792$ four-cycles.

\begin{proposition}\label{prop:q8struct}
The $680$-edge $C_4$-free subgraph of $Q_8$ has the following properties:
\begin{enumerate}[label=(\roman*)]
  \item Degree sequence $\{4^8,\,5^{160},\,6^{88}\}$; degree sum $1360=2\cdot680$.
  \item Bipartite, with $\lambda_1+\lambda_n=0$.
  \item Local maximality: every non-edge of $Q_8$ creates a $C_4$ upon addition.
  \item Edge density $680/1024=66.4\%$, compared to $304/448=67.9\%$
    for the $Q_7$ construction.
\end{enumerate}
\end{proposition}

\begin{remark}
Since $8$ is not a power of $4$, the exact BHN formula does not apply for
$n=8$.  The applicable estimate is
$\tfrac{1}{2}(8+0.9\sqrt{8})\cdot128\approx674.9$,
which our 680-edge construction surpasses by 5.1 edges.
\end{remark}

\subsection{Local optimality of the 680-edge construction}

For the $680$-edge solution, every non-edge of $Q_8$ creates at least
one $C_4$ upon addition.  Among the $344$ non-edges:
exactly $1$ creates exactly $1$ new $C_4$; exactly $1$ creates exactly $2$;
the remaining $342$ create $3$ or more.

Over $1\,076$ independent simulated-annealing trials targeting $681$ edges,
the minimum $C_4$ violation count observed was $1$ (never $0$),
where the violation count denotes the number of four-cycles fully contained
in the candidate edge set.

\begin{conjecture}\label{conj:q8}
$\ex(Q_8,C_4)=680$.
\end{conjecture}

Table~\ref{tab:compare} compares the $Q_7$ and $Q_8$ constructions.

\begin{table}[ht]
\centering
\caption{Structural comparison of the $Q_7$ and $Q_8$ constructions.}
\label{tab:compare}\medskip
\begin{tabular}{lcc}
\toprule
Parameter & $Q_7$ (304 edges) & $Q_8$ (680 edges)\\
\midrule
Total edges $|E(Q_n)|$ & 448 & 1024\\
Achieved edges & 304 & 680\\
Edge density & 67.9\% & 66.4\%\\
Min degree & 4 & 4\\
Max degree & 5 & 6\\
Degree sequence & $\{4^{32},5^{96}\}$ & $\{4^8,5^{160},6^{88}\}$\\
BHN estimate & $\approx300.2$ & $\approx674.9$\\
Improvement & $+3.8$ & $+5.1$\\
\bottomrule
\end{tabular}
\end{table}

%% ============================================================
\section{Structural Classification of $Q_7$ Solutions}\label{sec:structure}
%% ============================================================

\subsection{Dimension profiles and the twenty types}

\begin{definition}[Dimension profile]
For $G\subseteq Q_7$, let $e_i(G)=|\{uv\in E(G):u\oplus v=2^i\}|$.
The \emph{dimension profile} is the sorted tuple
$\pi(G)=(e_{\sigma(0)}\ge\cdots\ge e_{\sigma(6)})$.
\end{definition}

The dimension profile is invariant under coordinate permutations (a subgroup
of $\mathrm{Aut}(Q_7)$ of order $7!=5040$), giving a coarse but computable
classification.

\begin{proposition}\label{prop:twenty}
Among the $19\,866$ solutions, the dimension profile takes exactly
$\mathbf{20}$ distinct values.
\end{proposition}

Table~\ref{tab:types} lists all 20 types with frequencies and structural invariants.

\begin{table}[ht]
\centering\small
\caption{The 20 dimension-profile types of 304-edge $C_4$-free subgraphs of
  $Q_7$, sorted by decreasing frequency.
  $H$: Shannon entropy of normalised profile (bits).}
\label{tab:types}\medskip
\begin{tabular}{clrrl}
\toprule
Rank & Profile $\pi$ & Solutions & $H$ & $\lambda_1$ (mean)\\
\midrule
 1 & $[44,44,44,44,43,43,42]$ & 3155 & 2.8072 & 4.787\\
 2 & $[45,45,45,43,43,42,41]$ & 2913 & 2.8065 & 4.788\\
 3 & $[45,45,45,43,43,43,40]$ & 2200 & 2.8063 & 4.786\\
 4 & $[44,44,44,44,44,43,41]$ & 2116 & 2.8069 & 4.787\\
 5 & $[46,46,46,42,42,42,40]$ & 1756 & 2.8053 & 4.788\\
 6 & $[46,46,46,42,42,41,41]$ & 1181 & 2.8054 & 4.789\\
 7 & $[44,44,44,44,44,44,40]$ & 1085 & 2.8066 & 4.787\\
 8 & $[45,45,45,43,42,42,42]$ & 1064 & 2.8066 & 4.787\\
 9 & $[45,45,44,43,43,43,41]$ &  974 & 2.8067 & 4.787\\
10 & $[44,44,44,44,44,42,42]$ &  931 & 2.8070 & 4.787\\
11 & $[47,47,47,41,41,41,40]$ &  902 & 2.8037 & 4.789\\
12 & $[45,45,43,43,43,43,42]$ &  439 & 2.8069 & 4.787\\
13 & $[45,44,44,43,43,43,42]$ &  307 & 2.8070 & 4.787\\
14 & $[46,45,45,42,42,42,42]$ &  236 & 2.8063 & 4.787\\
15 & $[44,44,44,43,43,43,43]$ &  163 & 2.8073 & 4.787\\
16 & $[47,47,44,43,41,41,41]$ &  153 & 2.8050 & 4.788\\
17 & $[46,46,44,42,42,42,42]$ &  107 & 2.8062 & 4.787\\
18 & $[48,48,48,40,40,40,40]$ &  101 & 2.8014 & 4.789\\
19 & $[46,46,43,43,42,42,42]$ &   44 & 2.8063 & 4.787\\
20 & $[46,46,45,42,42,42,41]$ &   39 & 2.8058 & 4.787\\
\midrule
   & \textbf{Total}           & \textbf{19866}\\
\bottomrule
\end{tabular}
\end{table}

Observations: (1) $\lambda_1$ lies in $[4.786,4.789]$ across all 20 types.
(2) $H$ ranges from $2.801$ to $2.807$, all near $\log_27\approx2.807$.
(3) Type~18 ($[48,48,48,40,40,40,40]$) has 3-fold symmetry.

\subsection{Shared structural properties}

\begin{proposition}\label{prop:rigid}
Every solution in the $19\,866$ satisfies:
\begin{enumerate}[label=(\roman*)]
  \item Degree sequence $\{4^{32},5^{96}\}$.
  \item $\lambda_1+\lambda_n=0$ (bipartite), $\lambda_1\approx4.787$.
  \item $\mathrm{Tr}(A^4)=5216$ ($C_4$-free trace equality).
  \item Local maximality: no edge of $Q_7$ can be added without creating a $C_4$.
  \item Every edge of $Q_7$ appears in at least one solution
    (verified by taking the union of the edge sets over all $19\,866$ solutions).
\end{enumerate}
\end{proposition}

\subsection{Solution landscape}\label{sec:landscape}

Pairwise Hamming distances $d_H(G,G')=|E(G)\triangle E(G')|$
over a random sample of $5\,000$ pairs range from $36$ to $260$
(mean $\approx195.4$, median $196$).

Over $1\,076$ independent trials targeting $305$ edges,
the minimum $C_4$ violation count was $2$ (never $0$).

\begin{conjecture}\label{conj:q7}
$\ex(Q_7,C_4)=304$.
\end{conjecture}

%% ============================================================
\section{Computational Method}\label{sec:method}
%% ============================================================

\subsection{Two-phase simulated annealing}

\noindent\textbf{Phase 1 (Penalty SA).}
We minimise $f(E)=-|E|+\lambda V$ where $V$ is the number of $C_4$ violations
and $\lambda\in[0.30,0.90]$.
Temperature schedule: geometric from $T_0\in[0.20,4.00]$ to
$T_1\in[0.001,0.030]$ over $3$--$15$ million steps.
A move toggles a single edge; the change in $f$ is computed incrementally
in $O(\deg)$ time using precomputed $C_4$-to-edge incidence lists.

\noindent\textbf{Phase 2 (Swap SA).}
Edge count is fixed; swap moves (remove one edge, add one non-edge)
minimise $V$.
Parameters: up to $2\times10^8$ steps for $Q_7$ at 305 edges;
$5\times10^7$ for $Q_8$ at 681 edges.

\noindent\textbf{Diversification.}
Before each trial, a random element of $\mathrm{Aut}(Q_n)$
(random bit permutation composed with random bit flip) is applied to
the current best solution.  This prevents repeated exploration of the
same automorphism class.

\subsection{Comparison with known lower bounds}

\begin{align*}
n=7:&\quad\tfrac12(7+0.9\sqrt7)\cdot64\approx300.2,\quad\text{our bound: }304\;(\Delta=+3.8,\,+1.25\%),\\
n=8:&\quad\tfrac12(8+0.9\sqrt8)\cdot128\approx674.9,\quad\text{our bound: }680\;(\Delta=+5.1,\,+0.75\%).
\end{align*}

\subsection{Regularity trend across dimensions}

Degree sequences across dimensions suggest increasing regularity:
$Q_6$: $\{4^{32},5^{32}\}$;
$Q_7$: $\{4^{32},5^{96}\}$;
$Q_8$: $\{4^8,5^{160},6^{88}\}$.
The degree support $\{n-3,n-2,n-1\}$ may reflect an underlying structural
principle for optimal $C_4$-free subgraphs of $Q_n$.

%% ============================================================
\section{Computational Proof that $\ex(Q_6,C_4)=132$}\label{sec:q6}
%% ============================================================

Harborth and Nienborg~\cite{HN94} proved $\ex(Q_6,C_4)=132$ by combinatorial
arguments.  We reproduce this result computationally, providing an independent
verification and a template for small hypercubes.

\subsection{Lower bound: explicit construction}

Simulated annealing (Section~\ref{sec:method}) finds a $C_4$-free subgraph of
$Q_6$ on 132 edges, certified by exhaustive enumeration of all 240 four-cycles.
The explicit edge list is provided in \texttt{q6\_edges\_132.jsonl}.

\subsection{Upper bound: integer linear program}

Label the 192 edges of $Q_6$ as $x_1,\ldots,x_{192}\in\{0,1\}$.
For each of the 240 four-cycles $\{e_1,e_2,e_3,e_4\}$,
the constraint $x_{e_1}+x_{e_2}+x_{e_3}+x_{e_4}\le3$ ensures $C_4$-freeness.
The ILP is:
\[
  \max\sum_{i=1}^{192}x_i \quad\text{subject to}\quad
  \sum_{e\in C}x_e\le3\;\;\forall C\in\mathcal{C}_6,\quad x_i\in\{0,1\}.
\]
This formulation has 192 binary variables and 240 constraints.
The MPS file \texttt{q6\_ilp.mps} is provided for independent verification.

\begin{proposition}\label{prop:q6ilp}
The optimal value of the above ILP is $132$, establishing $\ex(Q_6,C_4)\le132$.
Combined with the lower bound construction, $\ex(Q_6,C_4)=132$.
\end{proposition}

On standard desktop hardware, standard MIP solvers such as SCIP, HiGHS,
or Gurobi verify this instance quickly.

%% ============================================================
\section{Open Problems}\label{sec:open}
%% ============================================================

\begin{enumerate}
  \item Prove or disprove $\ex(Q_7,C_4)=304$ (Conjecture~\ref{conj:q7}).
  \item Determine $\ex(Q_8,C_4)$ exactly (Conjecture~\ref{conj:q8}).
  \item What is the number of $\mathrm{Aut}(Q_7)$-orbits among all 304-edge
    $C_4$-free subgraphs of $Q_7$?
  \item Is the degree sequence $\{4^{32},5^{96}\}$ necessary for any locally
    maximal $C_4$-free subgraph of $Q_7$ on 304 edges?
  \item Can the flag algebra method of~\cite{Bab12,BHLL12} yield tight bounds
    for small $n\le8$?
  \item Does the degree sequence pattern $\{(n-3)^a,(n-2)^b,(n-1)^c\}$
    continue for $n=9,10,\ldots$?
\end{enumerate}

\section*{Acknowledgements}
The author thanks B.~Lidick\'y (Iowa State University) for the endorsement
and for pointing out the reference~\cite{Bab12}.
The author also thanks the open-source community for the HiGHS and SCIP
solvers used in ILP experiments.

\begin{thebibliography}{99}
\bibitem{Bab12}
R.~Baber,
Tur\'an densities of hypercubes,
arXiv:1201.3587, 2012.

\bibitem{BHLL12} J.~Balogh, P.~Hu, B.~Lidick\'y, and H.~Liu, Upper bounds on
  the size of 4- and 6-cycle-free subgraphs of the hypercube,
  \textit{Eur.\ J.\ Combin.}\ \textbf{35} (2014), 75--85.
\bibitem{BHN95} P.~Brass, H.~Harborth, and H.~Nienborg, On the maximum number
  of edges in a $C_4$-free subgraph of $Q_n$,
  \textit{J.\ Graph Theory}\ \textbf{19} (1995), 17--23.
\bibitem{Chu92} F.~Chung, Subgraphs of a hypercube containing no small even
  cycles, \textit{J.\ Graph Theory}\ \textbf{16} (1992), 273--286.
\bibitem{EMR92} M.\,R.~Emamy-K., P.~Guan, and P.\,I.~Rivera-Vega, On the
  characterization of the maximum squareless subgraphs of the 5-cube,
  \textit{Congr.\ Numer.}\ \textbf{88} (1992), 97--109.
\bibitem{Erd84} P.~Erd\H{o}s, On some problems in graph theory, combinatorial
  analysis and combinatorial number theory, in: \textit{Graph Theory and
  Combinatorics} (B.~Bollob\'as, ed.), Academic Press, 1984, pp.~1--17.
\bibitem{HN94} H.~Harborth and H.~Nienborg, Maximum number of edges in a
  six-cube without four-cycles, \textit{Bull.\ ICA}\ \textbf{12} (1994), 55--60.
\bibitem{TW09} A.~Thomason and P.~Wagner, Bounding the size of square-free
  subgraphs of the hypercube, \textit{Discrete Math.}\ \textbf{309} (2009),
  1730--1735.
\end{thebibliography}
\end{document}
